\section{An explicit Pless neighbour and its minimal shell}\label{app:pless-neighbour}

This appendix establishes the minimum and shell size in an explicit
integer model of an even neighbour. The theta-series argument applies
to every extremal even unimodular lattice of rank 48, including the
catalogue lattice in Section~\ref{app:p48} and the QR neighbour in
Section~\ref{sec:projection45}.

\subsection{The specified even neighbour}
Index the rows and columns of a $24\times24$ matrix $S$ by
$\{\infty\}\cup\F_{23}$, in that order. With $\chi$ the quadratic
character modulo $23$, set
\begin{equation}\label{eq:paley}
 S_{\infty\infty}=0,\quad S_{\infty j}=1,\quad S_{i\infty}=-1,
 \quad S_{ij}=\chi(i-j)\quad(i,j\in\F_{23}).
\end{equation}
The identity $SS^{\mathsf T}=23I$ follows from the quadratic-character
sum $\sum_x\chi(i-x)\chi(j-x)=-1$ for $i\ne j$. Consequently the
ternary code
\[
 C_3=\{(u,uS):u\in\F_3^{24}\}\subset\F_3^{48}
\]
is self-dual: its systematic generator has rank $24$, and
$I+SS^{\mathsf T}=24I=0$ over $\F_3$. Put
\begin{equation}\label{eq:construction-A}
 M=\{z\in\Z^{48}:z\bmod3\in C_3\},\qquad L=M/\sqrt3.
\end{equation}
The row basis
\[
 \begin{pmatrix}I_{24}&S\\0&3I_{24}\end{pmatrix}
\]
has determinant $3^{24}$. The code orthogonality makes $L$ integral,
and the scaling gives covolume one. It is odd, since it contains
$3e_0/\sqrt3$ of squared norm three.

Let $L_0=\{x\in L:x\cdot x\in2\Z\}$ and define
\begin{equation}\label{eq:characteristic}
 c=(1^{24},-23,1^{23}),\qquad g=c+6e_0,
 \qquad L'=L_0\cup\left(\frac{g}{2\sqrt3}+L_0\right).
\end{equation}
The characteristic congruence
$c\cdot z/3\equiv z\cdot z/3\pmod2$ for $z\in M$ is checked on
the displayed basis. This suffices since the squared norm modulo two
is additive on an integral lattice. Also $g/\sqrt3\in L_0$,
$g\cdot g/3=208$, and $g/(2\sqrt3)$ has squared norm $52$.
Its product with $z/\sqrt3\in L_0$ is integral by the characteristic
congruence. Thus~\eqref{eq:characteristic} is an even lattice, obtained
by extending the index-two sublattice $L_0$ by two cosets, and has
covolume one.

The matrix~\eqref{eq:paley} is the defining matrix of the Pless symmetry
code $C(23)$, with the same signs and coordinate order as
\cite[Section~1]{tonchev}. In particular, $C_3$ has minimum Hamming
weight $15$~\cite{pless,tonchev}. The sum of the rows of $(I\mid S)$
is $c$, so $C_3$ contains the all-one word. We shall also use the
following consequence of the complete weight enumerator of an extremal
ternary self-dual $[48,24,15]$ code containing that word:
\begin{equation}\label{eq:binary-words}
 \sum_{b\in C_3\cap\{0,1\}^{48}} t^{\wt(b)}
   =1+94t^{24}+t^{48}.
\end{equation}
This is the two-symbol specialization of
\cite[equation~(20)]{munemasa-tamura}. These two code-theoretic facts
determine which of the possible short vectors can occur in the even
neighbour.

\begin{theorem}[The mother shell]\label{thm:p48-shell}
The lattice $L'$ in~\eqref{eq:characteristic} is even unimodular,
has minimum squared norm $6$, and has exactly $52416000$ vectors of
squared norm $6$.
\end{theorem}
\begin{proof}
Evenness and unimodularity were established above. Suppose first that
$z/\sqrt3\in L_0$ has squared norm $2$ or $4$. Then $z\cdot z$ is
$6$ or $12$. A nonzero residue $z\bmod3$ would have Hamming weight
at least $15$, whereas its weight is at most $z\cdot z$. Hence
$z=3w$ for some integer vector $w$. Its squared norm is $3w\cdot w$,
which is even only when $w\cdot w$ is even, and is therefore at least
$6$ if nonzero.

A vector in the other coset has the form $y/(2\sqrt3)$, where
$y=g+2z$ and $z/\sqrt3\in L_0$. All $48$ coordinates of $y$ are odd,
so its squared norm is at least $4$. Equality would force
$y\in\{1,-1\}^{48}$. Let $b$ be the incidence vector of its negative
coordinates and put $w=\wt(b)$. Since $g\equiv\boldsymbol1\pmod3$,
the identity $z=(y-g)/2$ gives $-z\equiv b\pmod3$. Consequently
$b\in C_3$, and~\eqref{eq:binary-words} implies $w\in\{0,24,48\}$.
On the other hand, $c_i\equiv1\pmod6$, $\sum_i c_i=24$, and
$c\cdot g=582$. It follows that
\begin{equation}\label{eq:odd-parity}
 c\cdot z
 =\frac{\sum_i c_i-c\cdot g}{2}-\sum_{i:b_i=1}c_i
 \equiv3-w\equiv3\pmod6.
\end{equation}
This contradicts $z/\sqrt3\in L_0$, which requires
$c\cdot z\equiv0\pmod6$. Thus neither coset contains a vector of
squared norm $2$ or $4$. The vector $\sqrt3(e_0+e_1)\in L_0$
has squared norm $6$, proving the asserted minimum.

It remains to count the minimum shell. Write
$\Theta_{L'}(q)=\sum_{x\in L'}q^{x\cdot x/2}$. The theta series of
an even unimodular lattice of rank $48$ is a modular form of weight
$24$~\cite{conway1999sphere,mos}. The space has basis
$E_4^6,E_4^3\Delta,\Delta^2$. Its constant term and the vanishing
coefficients of $q$ and $q^2$ therefore give
\begin{equation}\label{eq:theta48}
 \Theta_{L'}=E_4^6-1440E_4^3\Delta+125280\Delta^2.
\end{equation}
Substituting
$E_4=1+240q+2160q^2+6720q^3+O(q^4)$ and
$\Delta=q-24q^2+252q^3+O(q^4)$ yields
\[
 \Theta_{L'}=1+52416000q^3+O(q^4),
\]
as required.
\end{proof}

Within this Pless-code model, the shift can equivalently be written
using $g_*=\boldsymbol1+6e_0$: indeed,
$(g-g_*)/2=-12e_{24}\in\sqrt3L_0$, so $g$ and $g_*$ define the
same coset of $L_0$. This is the all-one neighbour construction in
\cite{munemasa-tamura}; the parity calculation above explains directly
why this choice excludes the norm-four vectors. The liftings of
Section~\ref{app:p48} use the catalogue Gram basis and its published
minimum directly; the projection of Section~\ref{sec:projection45}
uses the separately specified QR neighbour.

The norm-six shell of $L'$, divided by $\sqrt6$,
is a kissing configuration $C$ of size $52416000$. Indeed, for two
distinct shell vectors $x,x'$, the nonzero lattice vector $x-x'$
has squared norm at least six, so $x\cdot x'\le3$.
